\section{Instrument detail}
\label{app:instrdetail}

The axis definitions are in Table~\ref{tab:axes}. The full operational judging
prompts and historical anchor rules are preserved in Appendix~\ref{app:axisprompts}.


\section{Positioning, and the mechanisms bounded}
\label{app:positioning}

\begin{table*}[t]
\centering\small
\begin{tabularx}{\textwidth}{@{}l l X@{}}
\toprule
Work & What it manipulates & Specific assumption this paper does not share \\
\midrule
\citet{sneakyprompt} & prompt tokens (RL search) & Optimises the lexical surface; a token substitution has a bounded search space and is reportable as a rate. \\
\citet{mmadiffusion}, \citet{ringabell} & weights or concept scores & White-box. We hold no gradients, weights or scores. \\
\citet{janus} & distribution, image filter folded in & Our evaluation additionally reports which material, exposure, and pose requirements the output meets. \\
\citet{universalperturb}, \citet{clipproxy} & pixels, directly & Require pixel control. \S\ref{app:noise} measures what remains when the attacker may only \emph{ask} for the perturbation, in this rendering style. \\
\citet{aestheticcamouflage}, \citet{etch} & semantics, style & Nearest prompt-level relatives; neither manipulates the language description of an \emph{imaging artefact} as carrier, nor reports how the content produced is scored. \\
\citet{visualjailbreak} & editing models, 15 categories & Genuinely multi-dimensional, but the dimensions are \emph{labels for the output}, not separately assessed material, exposure, and pose requirements. \\
\midrule
\textbf{this paper} & the prompt's referential frame & No optimisation, no gradients, no pixels. We contribute the attack, the decoupled evaluation of target attainment, and a measured boundary on the two alternatives the literature reaches for first. \\
\bottomrule
\end{tabularx}
\caption{\textbf{Positioning.} The method and evaluation distinguish image
generation from the joint attainment of specified visual attributes.}
\label{tab:positioning}
\end{table*}

\section{Measurement detail}
\label{app:measdetail}

\section{Output profiles and batch records}
\label{app:anticorrdetail}

Batch-level records are provided in Appendix~\ref{app:ladder}; historical
per-image judging outputs are in Appendix~\ref{app:instances}. These records
support the distinction between image return and target attainment introduced
in \S\ref{sec:evalmotivation}.

\section{The length control}
\label{app:lengthcontrol}

\begin{table}[t]
\centering\small
\begin{tabular}{lrll}
\toprule
Cell & Chars & Content & Result \\
\midrule
\texttt{A\_short}          & 101     & none    & \textbf{3/3 pass} (30.6/21.2/18.8\,s) \\
\texttt{B\_long\_benign}   & 4{,}561 & none    & \textbf{3/3 pass} (30.1/26.8/29.0\,s) \\
\texttt{C\_short\_content} & 629     & present & blocked (75.2/75.3\,s) \\
\texttt{D\_long\_content}  & 4{,}561 & present & blocked \\
\bottomrule
\end{tabular}
\caption{\textbf{Length is not the gate; content is.} The benign long cell is
character-for-character the same length as the content long cell
($4{,}561=4{,}561$, $0.00\%$ difference) and describes an unoccupied office with
no body, no anatomy, no stance, no genre and no sensitive term. It passes $3/3$
while the content cell blocks. The short content cell is one seventh the length
of the control and still blocks.}
\label{tab:length}
\end{table}

\begin{figure}[t]
\centering
\includegraphics[width=\linewidth]{figures/figM_anchors.png}
\caption{\textbf{Two failures of our own instrument, quantified.}
(a) Anchor attendance on the 98-image re-judged corpus. \textbf{Four of the
twelve anchors never, or almost never, fire}---most of the rule is unexercised
on this corpus, which must be stated or a reader will take ``not observed'' for
``does not occur''. (b) Displacement of the material reading under three
observation arms on the same 20 images. Adding a mark moves the reading on 5;
\emph{cropping the carrier away} moves it on 12. \textbf{Part of the material
reading is supplied by the presentation rather than by the body}, showing why
the depicted figure must be assessed separately from its presentation.}
\label{fig:anchors}
\end{figure}

\begin{figure}[t]
\centering
\includegraphics[width=\linewidth]{figures/figO_judging.png}
\caption{\textbf{Configuration effects that masquerade as findings.} (a) Our own
first judging run used too small an output budget; the evidence field was
truncated and 19 images were recorded as null, which we briefly read as model
disagreement. With a larger budget the same batch returned the top level on 6/6
and upper-range levels rose from 6 to 35. \textbf{Every missing field is a parse
failure until proven otherwise.} (b) Repeat stability: five interleaved repeats
give zero level and zero anchor flips, while duplicates separated in time across
six rounds give 1 of 7 level flips and 3 anchor flips---\textbf{the time-dependent
component concentrates in anchors that require interpretation rather than
observation}.}
\label{fig:judging}
\end{figure}

\begin{figure}[t]
\centering
\includegraphics[width=\linewidth]{figures/figQ_threat.png}
\caption{\textbf{What the attacker has, and what the measurement can therefore
be.} The attack is strictly black-box and text-only. \textbf{The attackable
object is the response, not the decision behind it}, which is why this paper
scores the image the endpoint returns rather than attempting to attribute a
refusal to any particular filter.}
\label{fig:threat}
\end{figure}

\begin{figure*}[t]
\centering
\includegraphics[width=\textwidth]{figures/figE_prompt_anatomy.png}
\caption{\textbf{Prompt anatomy, and where each ablation acts.} Every prompt is
assembled from six segments. Each sweep in this paper varies exactly one segment
with the other five held byte-identical, which is what makes the comparisons
single-variable. The three brackets below the segments mark the three independent
ablations reported in \S\ref{app:ablationdetail}: the opening noun of the shell, the
pose clause (75 distinct clauses tested), and the number of stacked texture layers
in the carrier segment (ten layers are indistinguishable from one).}
\label{fig:anatomy}
\end{figure*}

\begin{figure*}[t]
\centering
\includegraphics[width=0.94\textwidth]{figures/figF_outputs.png}
\caption{\textbf{Outputs, by carrier.} Four passing images from the carrier sweep
of Appendix~\ref{app:ablationdetail}, each retaining the full figure and reaching the
exposure levels the instrument records. The carrier acts only outside the figure
silhouette in every case; the figure itself is explicitly exempted by the recipe
(\S\ref{app:noise}). These are outputs the endpoint returned with HTTP 200, shown
at the resolution it returned.}
\label{fig:outputs}
\end{figure*}

\begin{table}[t]
\centering\small
\begin{tabular}{@{}p{0.30\linewidth}p{0.62\linewidth}@{}}
\toprule
Extent & Claims in this paper \\
\midrule
\textbf{observed} &
\fTotal{} logged requests with HTTP state, latency and error text; \fFiles{}
result files across \fBatches{} batch directories; the retained output images; the
prompt files with byte-level identity assertions for every ablation pair; the
four-state census; the latency distributions; the noise calibration; and the
out-of-band counter control. \\
\midrule
\textbf{not verified} &
human construct validity; any provider's policy determination; the material axis
(presentation confound); the top-tier round's three-axis judgement; the cause of
the non-stationarity; and any statement about non-sexual NSFW categories. \\
\bottomrule
\end{tabular}
\caption{\textbf{What was measured, and what was not.} The two rows are disjoint;
nothing in the second row is used to support a claim in the first.}
\label{tab:extent}
\end{table}


\section{Target-control records}
\label{app:targetrecords}

The main-text examples use the saved outputs
\texttt{runs151/R11\_r1.png}, \texttt{runs151/R13\_r3.png}, and
\texttt{runs152/K34b\_r4.png}. The observations in
Table~\ref{tab:targetoutputs} describe the images below. They retain the
visible rendering style and body configuration as separate properties.
Geometric chest marks are recorded explicitly.

\begin{figure*}[t]
\centering
\includegraphics[width=0.31\textwidth]{figures/target_standing.png}\hfill
\includegraphics[width=0.31\textwidth]{figures/target_crouching.png}\hfill
\includegraphics[width=0.31\textwidth]{figures/target_rear.png}
\caption{\textbf{Original target-control outputs.} Left: batch 151, R11,
repeat 1. Centre: batch 151, R13, repeat 3. Right: batch 152, K34b,
repeat 4. These are the original images used for the attribute observations
in Table~\ref{tab:targetoutputs}.}
\label{fig:targetrecords}
\end{figure*}

Batch 156 contains 39 saved request records. Groups M, N, P, and G each
contain three pose settings with three attempts per setting. Group Q contains
three further pose-wording attempts. N11 and P11 each return one image;
the other groups return none. Table~\ref{tab:rendercontext} uses groups
M, N, P, and G. The source filenames and record hashes are frozen in
\texttt{restructure\_evidence.json}. These records supplement the earlier
corpus census rather than changing its frozen totals.

\section{Supplementary component ablations}
\label{app:ablationdetail}

The historical component sweeps vary the opening context, texture depth,
and carrier form separately. Their original summary is reproduced in
Figure~\ref{fig:suppcomponents}. These measurements concern returned images;
target-attribute attainment requires the separate checks in
\S\ref{sec:targetmatch}.

\begin{figure}[t]
\centering
\includegraphics[width=\linewidth]{figures/figD_three_arms.png}
\caption{\textbf{Component sweeps.} The panels vary the opening context,
texture-layer count, and carrier form. Each panel represents its own batch;
the layer sweep covers one through ten layers.}
\label{fig:suppcomponents}
\end{figure}


\section{Method showcase records}
\label{app:methodoutputs}

Figure~\ref{fig:methodcontrol} reproduces three saved generation outputs:
(a) \texttt{runs69/K\_seated\_lean\_r3.png},
(b) \texttt{runs81/S2\_r2.png}, and
(c) \texttt{runs74/lie\_one\_knee\_up\_r4.png}.
The prior image-level assessments in \texttt{gallery/\_expo\_judge.json}
record exposure level 6 for all three; the curated collection records pose
level 2. Reinspection of the original images confirms a modern drawn central
figure, identifiable nipple/areola and genital detail, and a sexualized posture
without visible sexual activity. Panel labels describe the visible configurations.
The source paths, file hashes, and per-image readings are preserved in
\texttt{method\_showcase\_evidence.json}.
